\section{Conclusion and Limitations}

Smart Coder's main contribution is not a claim of fully solved personalization; it is a governed adaptive-autonomy framework that makes oversight behavior explicit, persistent, and observable. By separating task contracts from softer preferences and by logging both model-initiated and policy-initiated check-ins, the system creates a concrete platform for studying how coding-agent autonomy should adapt over time rather than treating that behavior as an opaque model property.

The current prototype remains a heuristic baseline. Its approval policy uses manually chosen priors rather than learned calibration, its spec support is lightweight rather than fully structured, and it does not yet implement richer interrupt taxonomy, post-hoc correction, rewind, or asynchronous delegation. These are meaningful limitations, but they are also appropriate for a demo paper whose goal is to present a working artifact and an instrumented workflow rather than a finalized learning system.

The immediate next step is a controlled lab study using the packaged CLI, the demo repository, and exported traces from repeated sessions. That study will let us contribute a clearer taxonomy of user preference types, test whether different oversight styles emerge across users and tasks, and evaluate whether learned policies or fine-tuned models can outperform the current heuristic runtime. Beyond coding, the same framework may generalize to other agentic workflows where user preferences mix hard governance rules with softer behavioral conditioning.
